\section{6DOF Simulation}

In \texttt{C:\textbackslash 02-Dinamica6DOF}, the missile data, CG/inertia, and launch-condition files were edited. The script \texttt{VooUnico.m} was also created. The Simulink model \texttt{Dinamica6DOF.slx} was kept as provided. The \texttt{M\_aed.mat} matrix generated from DATCOM was copied to this folder before the simulation.

The test was run during the initial acceleration segment, with an initial speed of 10 m/s and a control deflection of \(\Delta Y = 2^\circ\). This low initial speed is a numerical test condition, not a claim about the actual launch speed. It lets the model start smoothly from the beginning of the booster phase and keeps the run focused on the consistency between the generated aerodynamic table and the 6DOF model. Aerodynamic forces in the first instants are therefore lower than they would be in an operational post-launch condition. The final time was limited to 2.8 s to keep Mach inside the main range of the aerodynamic table. At the end of the simulation, the missile reached Mach 0.876 and 294.9 m/s.

\begin{lstlisting}[language=Matlab]
C.Tfim = 2.8;
C.DeltaY = 2*D2R;
C.DeltaZ = 0*D2R;
C.Vmx0 = 10;
\end{lstlisting}

\subsection*{6DOF results}

\begin{figure}[H]
\centering
\includegraphics[width=0.9\textwidth]{figures/dof_alfa_delta.png}
\caption{Alpha and beta response to the applied control deflection.}
\end{figure}

\begin{figure}[H]
\centering
\includegraphics[width=0.86\textwidth]{figures/dof_alpha_delta_overlay.png}
\caption{Direct comparison between \(\alpha(t)\) and the applied \(\Delta Y(t)\).}
\end{figure}

\begin{figure}[H]
\centering
\includegraphics[width=0.9\textwidth]{figures/dof_velocidade_aceleracao.png}
\caption{Velocity, Mach number, and accelerations in the 6DOF test.}
\end{figure}

\begin{figure}[H]
\centering
\includegraphics[width=0.9\textwidth]{figures/dof_propulsao_massa.png}
\caption{Thrust, mass, CG, and inertia during the initial burn.}
\end{figure}
